\pagebreak
\nopagecolor
\begin{center}
  {\fontsize{30pt}{26pt}\selectfont 证明手册}
\end{center}
\separator
\section*{I. 代数}

\begin{questions}

\question{证明二次公式：若 $ax^{2}+bx+c=0$ ($a \neq 0$)，则 $x = \frac{-b \pm \sqrt{b^{2}-4ac}}{2a}$。}
\begin{solution}
    将方程两边同时除以 $a$：$x^{2} + \frac{b}{a}x + \frac{c}{a} = 0$。\\
    移项得：$x^{2} + \frac{b}{a}x = -\frac{c}{a}$。\\
    配方：在两边同时加上 $(\frac{b}{2a})^{2}$，得 $x^{2} + \frac{b}{a}x + (\frac{b}{2a})^{2} = \frac{b^{2}}{4a^{2}} - \frac{c}{a}$。\\
    整理为完全平方式：$(x + \frac{b}{2a})^{2} = \frac{b^{2}-4ac}{4a^{2}}$。\\
    两边开平方：$x + \frac{b}{2a} = \pm \frac{\sqrt{b^{2}-4ac}}{2a}$。\\
    最终得：$x = \frac{-b \pm \sqrt{b^{2}-4ac}}{2a}$。
\end{solution}

\question{证明立方和与立方差公式：$a^{3} \pm b^{3} = (a \pm b)(a^{2} \mp ab + b^{2})$。}
\begin{solution}
    以立方和为例，展开右侧多项式：\\
    $(a + b)(a^{2} - ab + b^{2}) = a(a^{2} - ab + b^{2}) + b(a^{2} - ab + b^{2})$ \\
    $= (a^{3} - a^{2}b + ab^{2}) + (a^{2}b - ab^{2} + b^{3})$ \\
    项相互抵消后得：$= a^{3} + b^{3}$。同理可证立方差公式。
\end{solution}

\question{证明对数性质：$\log_{a}xy = \log_{a}x + \log_{a}y$。}
\begin{solution}
    设 $M = \log_{a}x$，$N = \log_{a}y$，则 $x = a^{M}$，$y = a^{N}$。\\
    因此 $xy = a^{M} \cdot a^{N} = a^{M+N}$。\\
    根据对数定义，$\log_{a}xy = M + N = \log_{a}x + \log_{a}y$。
\end{solution}

\question{证明对数换底公式：$\log_{a}x = \frac{\log_{b}x}{\log_{b}a}$。}
\begin{solution}
    设 $y = \log_{a}x$，则 $x = a^{y}$。\\
    两边同时取以 $b$ 为底的对数：$\log_{b}x = \log_{b}(a^{y})$。\\
    利用幂规则：$\log_{b}x = y \log_{b}a$。\\
    解出 $y$：$y = \frac{\log_{b}x}{\log_{b}a}$，即 $\log_{a}x = \frac{\log_{b}x}{\log_{b}a}$。
\end{solution}

\question{证明等差数列求和公式：$S_{n} = \frac{n}{2}[2a + (n-1)d]$。}
\begin{solution}
    将 $S_{n}$ 正序与倒序排列相加：\\
    $S_{n} = a + (a+d) + \dots + [a+(n-1)d]$ \\
    $S_{n} = [a+(n-1)d] + [a+(n-2)d] + \dots + a$ \\
    两式相加，对应项的和均为 $2a + (n-1)d$，共有 $n$ 项：\\
    $2S_{n} = n[2a + (n-1)d]$ \\
    故 $S_{n} = \frac{n}{2}[2a + (n-1)d]$。
\end{solution}

\question{证明等比数列求和公式：$S_{n} = \frac{a(1-r^{n})}{1-r}$ ($r \neq 1$)。}
\begin{solution}
    列出前 $n$ 项和：$S_{n} = a + ar + ar^{2} + \dots + ar^{n-1}$。\\
    两边乘以 $r$：$r S_{n} = ar + ar^{2} + \dots + ar^{n}$。\\
    两式相减：$S_{n} - r S_{n} = a - ar^{n}$。\\
    提取公因式：$S_{n}(1-r) = a(1-r^{n})$。\\
    由于 $r \neq 1$，得 $S_{n} = \frac{a(1-r^{n})}{1-r}$。
\end{solution}

\question{证明自然数平方和公式：$\sum_{k=1}^{n}k^{2} = \frac{n(n+1)(2n+1)}{6}$。}
\begin{solution}
    利用恒等式 $(k+1)^{3} - k^{3} = 3k^{2} + 3k + 1$。\\
    对 $k=1$ 到 $n$ 求和：\\
    $\sum_{k=1}^{n} [(k+1)^{3} - k^{3}] = 3\sum k^{2} + 3\sum k + \sum 1$ \\
    左侧为裂项相消：$(n+1)^{3} - 1^{3} = 3\sum k^{2} + 3\frac{n(n+1)}{2} + n$ \\
    整理方程，解得 $\sum k^{2} = \frac{2(n+1)^{3} - 3n(n+1) - 2(n+1)}{6} = \frac{n(n+1)(2n+1)}{6}$。
\end{solution}

\question{证明棣莫弗定理：$[r(\cos\theta + i \sin\theta)]^{n} = r^{n}(\cos n\theta + i \sin n\theta)$。}
\begin{solution}
    利用欧拉公式 $e^{i\theta} = \cos\theta + i \sin\theta$。\\
    左式 $= (r e^{i\theta})^{n} = r^{n} e^{in\theta}$。\\
    再次应用欧拉公式：$r^{n} e^{in\theta} = r^{n}(\cos n\theta + i \sin n\theta)$。\\
    （或使用数学归纳法结合三角恒等式证明）。
\end{solution}

\question{证明概率加法公式：$P(A \cup B) = P(A) + P(B) - P(A \cap B)$。}
\begin{solution}
    由集合论可知 $A \cup B = A \cup (B \setminus (A \cap B))$，且 $A$ 与 $B \setminus (A \cap B)$ 是互斥的。\\
    因此 $P(A \cup B) = P(A) + P(B \setminus (A \cap B))$。\\
    又因为 $B = (B \setminus (A \cap B)) \cup (A \cap B)$ 且互斥，所以 $P(B) = P(B \setminus (A \cap B)) + P(A \cap B)$。\\
    代入得 $P(A \cup B) = P(A) + [P(B) - P(A \cap B)]$。
\end{solution}
\end{questions}

% \section*{I. 代数}

% \begin{questions}
% \end{questions}